\chapterimage{lan_party.jpg}
\chapter{Capa 2: Comunicació de xarxa (LAN)}\label{sec:layer2}

\begin{minipage}{0.4\linewidth}
\begin{center}
\begin{bytefield}{16}
\bitbox{16}{Layer 7: Application} \\
\bitbox{16}{Layer 4: Transport} \\
\bitbox{16}{Layer 3: Internet} \\
\bitbox{16}{\color{color1} Layer 2: Network (LAN)} \\
\bitbox{16}{Layer 1: Physical} \\
\end{bytefield}
\end{center}
\end{minipage}
\begin{minipage}{0.6\linewidth}
\begin{center}
\includegraphics[width=\linewidth]{network_layer2.pdf}
\end{center}
\end{minipage}

\vspace{-0.75cm}

% \section{Layer~2 overview}
\subsection*{Capacitats}

Els protocols de capa~2 permeten que diversos dispositius intercanviïn missatges
i s'identifiquin entre ells dins d'una xarxa local (LAN (xarxa d'àrea local)).
% 
Amb els protocols de capa~2 no es pot contactar \textit{directament} amb dispositius de fora de la LAN.

La capa~2 abstreu les capes superiors (incloses les aplicacions d'usuari) dels detalls específics del maquinari connectat.
Això fa més fàcil (és a dir, més barat) escriure el codi una vegada i executar-lo a tot arreu.

Els missatges de capa~2 tenen una longitud limitada. La quantitat màxima de dades de payload (càrrega útil)
que es pot inserir en un protocol de capa~2 s'anomena unitat màxima de transmissió (MTU),
\eg, $1500$~bytes per a Ethernet.

Alguns protocols de capa~2 també poden oferir autenticació d'usuaris i
protecció criptogràfica;
\conceptRef{wi-fi}{Wi-Fi} n'és un exemple destacat,
tot i que les xarxes amb cable també poden fer servir, \eg, IEEE 802.1X.

\vspace{-0.3cm}
\begin{center}
\includegraphics[width=0.9\linewidth]{protocols_layer2.pdf}
\end{center}
\vspace{-0.5cm}

\subsection*{Protocols}
S'han definit i s'estan definint molts protocols dins de la capa~2. Els més importants avui dia són:\\[-0.65cm]
\begin{itemize}
\item \conceptRef{ethernet}{Ethernet} (IEEE 802.3), tractat a la \secref{sec:layer2:ethernet}.
\item \conceptRef{wi-fi}{Wi-Fi} (\eg, IEEE 802.11be, també conegut com a Wi-Fi 7, del 2024).
\end{itemize}

Dins de TCP/IP, la capa~2 encapsula les PDUs dels protocols següents
(tots dos tractats al capítol~\ref{sec:layer3}):\\[-0.65cm]
\begin{itemize}
  \item Datagrames \concept{IPv4} i \concept{IPv6}.
  \item Missatges ARP.\\[-0.75cm]
\end{itemize}

\vspace{0.5cm}

\begin{remark}
L'ordre \inlineCode{ip link} mostra les teves interfícies i permet configurar-les,
incloent-hi l'adreça \concept{MAC}.

El teu sistema operatiu normalment assigna automàticament noms interns a les interfícies segons el seu tipus,
\eg, \inlineCode{eth0} o \inlineCode{wlan1}.
% 
Normalment es poden canviar els noms de les interfícies, perquè aquests noms
no formen part de TCP/IP; només es fan servir dins d'aquell sistema operatiu.
\end{remark}


\section{Adreçament de capa~2}\label{sec:layer2:layer_2_addressing}
En la majoria de tecnologies LAN, les adreces de capa~2 són identificadors numèrics únics per a cada dispositiu dins de la LAN. 
El tipus més habitual són les adreces MAC d'Ethernet (tècnicament EUI-48), que ocupen $6$~bytes (\ie, $48$~bits).
També es fan servir fora d'Ethernet, \eg, a Wi-Fi i Bluetooth.

Les adreces MAC es presenten a les persones sovint en el format següent:\\\otherBase{01:23:45:67:89:ab}.
Algunes adreces tenen un significat especial i no són identificadors de dispositiu vàlids.
Per exemple, \otherBase{ff:ff:ff:ff:ff:ff} és l'adreça de \concept{broadcast} (difusió), que vol dir ``tothom a la LAN''. 
% 
A més, \otherBase{00:00:00:00:00:00} es fa servir per representar una adreça MAC desconeguda
% 
(consulta l'\href{https://www.iana.org/assignments/ethernet-numbers/ethernet-numbers.xml}{\underline{RFC~9542}} 
per a una descripció completa de les adreces MAC reservades, incloses les de \conceptRef{unicast/multicast}{multicast} (enviament a un grup)).

El maquinari de xarxa ve amb una adreça MAC integrada, única i predefinida, a punt per fer-se servir.
Les MAC dels dispositius es poden configurar a nivell del sistema operatiu, però es mantenen constants mentre s'utilitzen.


\section{Protocol de capa 2 Ethernet}\label{sec:layer2:ethernet}

\subsection{Format de paquet}\label{sec:layer2:ethernet:packet_format}
La capa~2 d'Ethernet defineix el format següent, que es fa servir en tots els missatges.
Els paquets d'aquest tipus s'anomenen frames (trames):\\[-0.25cm]

\begin{center}
\begin{bytefield}[bitheight=3em]{48}
\bitheader{0,7,8,15,16,23,24,31,32,39,40,47}\\
\bitbox{48}{Destination MAC Address} \\ 
\bitbox{48}{Source MAC Address} \\
\bitbox{16}{Payload type} & 
\bitbox[tlr]{32}[bitheight=3em]{Payload data\\{\scriptsize(variable length, byte aligned)}} \\
\bitbox[lbr]{40}{\hspace{7cm}\raisebox{0.25cm}{$\vdots$}} & \bitbox[lt]{8}{} \\
% 
% \bitbox[tr]{32}{Payload data (variable length)} \\
% \bitbox[lr]{48}{\raisebox{1em}{$\vdots$}} \\
% \bitbox[lbr]{24}{} & \bitbox[t]{24}{} 
\end{bytefield}
\end{center}

\begin{itemize}
\item \textbf{MAC de destinació}: l'adreça MAC del dispositiu de destinació, 
  o \otherBase{ff:ff:ff:ff:ff:ff} per a \concept{broadcast}.
\item \textbf{MAC d'origen}: la MAC del dispositiu que envia el frame.
\item \textbf{Tipus de payload}: identificador del protocol 
  encapsulat al payload,\\\eg, 
  \otherBase{0x0800}~$\rightarrow$~\concept{IPv4}, 
  \otherBase{0x86DD}~$\rightarrow$~\concept{IPv6} i
  \otherBase{0x0806}~$\rightarrow$~ARP.
  De vegades s'anomena \href{https://en.wikipedia.org/wiki/EtherType#Values}{\underline{EtherType}}.
  
\item \textbf{Payload}: contingut que es demana enviar, \eg, per part d'IPv4 o ARP.
  La longitud màxima d'aquest camp, la MTU, és de $1500$~bytes per a Ethernet.
\end{itemize}

\begin{exercise}\ \\[-0.5cm]
\begin{itemize}
\item Quantes adreces MAC diferents hi ha a Ethernet (incloent-hi els blocs reservats)?
\item N'hi ha prou perquè cada dispositiu amb connexió a internet de la Terra tingui una MAC única?
\end{itemize}
\end{exercise}

\begin{exercise} 
És vàlid el frame següent?
% 
\begin{center}
\otherBase{c2 21 90 15 \quad bc b6 42 40 \quad e5 f7 5e 32 \quad 08 42 77 57}\\
\otherBase{03 ca 27 6b \quad 1f 97 e9 02 \quad 24 7f 80 fa \quad 8b 61 59 2a}\\
\otherBase{81 b6 73 b0 \quad b4 e1 75 4f \quad ef 40 dd 9f \quad 34 3c 48 67}
\end{center}
\end{exercise}

\newpage
\subsection{Funcionament}\label{sec:layer2:ethernet:operation}

Els missatges dins d'una LAN Ethernet normalment s'envien en mode \conceptRef{unicast/multicast}{unicast} (enviament a un únic destinatari), \ie, són missatges d'un a un.
Això permet als switches (commutadors) minimitzar el consum d'energia i d'\conceptRef{bandwidth}{amplada de banda}.
% 
En \conceptRef{unicast/multicast}{unicast}, el node del dispositiu d'origen ha de conèixer la MAC de la destinació.
% 
\begin{center}
\includegraphics[width=\linewidth]{ethernet_message_unicast.eps}
\end{center}

Ethernet també admet \concept{broadcast} a tots els nodes de la LAN. En aquest cas, l'adreça MAC
de destinació ha de ser \otherBase{ff:ff:ff:ff:ff:ff}.
% 
Quan un dispositiu rep un frame, el descarta tret que el camp de MAC de destinació del frame
sigui idèntic a la seva pròpia adreça MAC, o de \concept{broadcast}. 

Un node també descartarà un frame amb un camp de tipus de payload (\conceptRef{ethertype}{EtherType}) no admès. 
% 
Si el tipus de payload és admès, es fa servir per decidir quina part del sistema operatiu rep el missatge,
\eg, la pila \concept{IPv4} o el subsistema ARP.

\begin{exercise} Continuant l'exemple de la figura anterior, la targeta de xarxa del node~B rep
un frame que comença amb els bytes següents. L'hauria d'acceptar B?
\begin{center}
\otherBase{00 11 22 00 \quad 00 0a ff ff \quad ff ff ff ff \quad 08 DD \ldots}
\end{center}
\end{exercise}

\begin{exercise}
Els scripts següents mostren com enviar i rebre frames \conceptRef{ethernet}{Ethernet} en brut. 
El client envia $5$ frames idèntics d'\conceptRef{ethertype}{EtherType} \otherBase{0x1234} i després acaba.
El servidor espera fins que rep $5$ frames Ethernet d'aquest tipus (comprovat a les línies 12-13) i després acaba.
% 
\begin{itemize}
\item Canvia l'adreça MAC i el nom d'interfície del client pels teus.
\item Fes que el servidor rebutgi els frames que no li són adreçats.
\item Amplia aquests scripts per implementar una aplicació senzilla amb capacitats LAN.\\
  Exemples suggerits:
    \begin{itemize}
    \item Un servei d'eco de capa~2 que distingeixi entre peticions i respostes.
    \item Un xat peer-to-peer (de parell a parell) de capa~2 que inclogui el sobrenom de l'origen en tots els missatges.
    \item Un servei de calculadora de capa~2 que admeti operacions aritmètiques bàsiques\\
      (\inlineCode{+}, \inlineCode{-},\inlineCode{*}, divisió entera \inlineCode{//} i, opcionalment, divisió \inlineCode{/}).
    \end{itemize}
\item Què falta perquè la teva aplicació es pugui connectar a la resta d'Internet més enllà de la teva LAN?
\end{itemize}

\label{ex:layer2:echo}
\end{exercise}

\vfill

\begin{center}
\showCode{snippets/ethernetclient.py}
\end{center}

\begin{center}
\showCode{snippets/ethernetserver.py}
\end{center}

\begin{remark}
Hauràs d'executar aquests scripts amb privilegis, \eg, amb \inlineCode{sudo}.
% 
Com a alternativa, pots afegir de manera permanent les capacitats \texttt{CAP\_NET\_RAW}
al teu binari de python amb
\begin{center}
\inlineCode{sudo setcap cap_net_raw+ep ./venv/bin/python}.  
\end{center}
% 
Després, pots executar \inlineCode{./venv/bin/python script.py} directament sense \inlineCode{sudo}.\\[-0.5cm]
\end{remark}
